% =====================================================================
% CONCLUSION
% =====================================================================

\section{Conclusion}
\label{sec:conclusion}

We studied which component of the test loss is responsible for the
double-descent peak. We gave an upper bound on the MLE excess test cross-entropy
that decomposes into four interpretable terms: an approximation-and-shift
term, a minimax batch regret, a prior cost, and an MLE-vs-Bayes gap.
Under the posterior-sampling assumption, we also lower-bounded the
MLE-vs-Bayes gap by the expected MLE-vs-ensemble loss difference in the
well-specified setting. Through experiments under
fully random labels, where the approximation-and-shift term is zero by
construction, we found a large SGD-vs-ensemble gap near interpolation,
supporting the MLE-vs-Bayes gap as a major contributor to the peak. For practitioners
working in the interpolation regime, this picture suggests that the
double-descent peak can depend strongly on the prediction rule, and that
probability-averaged ensembles can attenuate it.

\textbf{Limitations.} Our experimental sweeps are less extensive than
prior empirical studies of double descent \citep{nakkiran2021deep}, do
not include the largest modern architectures, and use eight SGLD chains
per CNN width. The gap lower bound assumes an ideal posterior sampler.
Our finite SGLD runs may retain mixing and discretization error;
posterior consistency is assumed rather than proved for the CNNs studied here.
The longer-run and smaller-step checks support robustness of the observed
comparison, but do not certify the measured finite-test-set gap as a lower
bound. The universal compression decomposition relies
on a supremum over the model family in the batch-regret and MLE-vs-Bayes
terms, resulting in an upper bound rather than an equality; this is the
price for cleanly separating the batch regret from the MLE-vs-Bayes gap.
We do not establish numerical tightness for deep networks or directly
estimate (II), (II'), or the supremum in (III). Finally, while the
random-labels experiment eliminates (I) and shows a large empirical gap,
it does not establish that (III) dominates across all settings.
